% NOTE (2026-09-19): the Lean files this blueprint points at are Cowork drafts
% that have not been compiled yet (see docs/TASKS.md, T1).  Therefore **no node
% carries \leanok**.  Add \leanok to a node when, and only when, the
% corresponding declaration compiles and `leanblueprint checkdecls` resolves it.

\chapter{The model and the propagator}

\begin{definition}[Block covariance matrix, \S2.1]\label{def:SB}
  \lean{RBM.SB}
  Let $L\ge 3$ and work on the discrete torus $\ZL$.  Set
  \[ \SB_{ab} \;=\; \tfrac15\,\mathbf 1\!\left(|a-b|_L\le 1\right),\qquad a,b\in\ZL, \]
  the transition matrix of the lazy nearest-neighbour random walk on $\ZL$.
  Equivalently, $\SB$ is the circulant matrix on the group $\ZL$ generated by the
  kernel supported on $\{(0,0),(\pm1,0),(0,\pm1)\}$ with value $1/5$.
\end{definition}

\begin{remark}\label{rem:L3}
  The paper leaves the hypothesis $L\ge 3$ implicit.  It is needed: for $L\le 2$
  the five lattice points above are not distinct and $\SB$ fails to be stochastic.
  See \texttt{docs/paper-deltas.md}, item 1.
\end{remark}

\begin{definition}[Periodic distance]\label{def:dist}
  \lean{RBM.zdist2}
  $|x|_L$ denotes the periodic distance on $\ZL$.  The paper uses the periodic
  Euclidean norm and notes that the periodic $L^1$ and $L^2$ norms differ by at
  most $\sqrt2$, which is absorbed into the constants; we formalize the $L^1$
  one.  See \texttt{docs/paper-deltas.md}, item 2.
\end{definition}

\begin{lemma}[Elementary properties of $\SB$]\label{lem:SB-basic}
  \lean{RBM.SB_isSymm, RBM.SB_apply_add_right, RBM.sum_SB_row, RBM.SB_mulVec_one,
    RBM.card_sbSupport}
  \uses{def:SB}
  $\SB$ is symmetric and translation invariant, each row sums to $1$, and
  $\SB\mathbf 1=\mathbf 1$.
\end{lemma}

\begin{definition}[The propagator, Definition \texttt{def\_Theta}]\label{def:Theta}
  \lean{RBM.Theta}\uses{def:SB}
  For $\xi\in\CC$ with $|\xi|<1$,
  \[ \Th^{(B)}_\xi := \frac{1}{1-\xi\,\SB}. \]
  In Lean this is \texttt{Ring.inverse} of $1-\xi\SB$ in the matrix ring with the
  $\ell^\infty$ operator norm; $\|\SB\|=1$ makes $1-\xi\SB$ a unit.
\end{definition}

\begin{lemma}[$\Th$ is the unique inverse]\label{lem:Theta-unique}
  \lean{RBM.Theta_mul, RBM.mul_Theta, RBM.eq_Theta_of_mul}\uses{def:Theta}
  $\Th^{(B)}_\xi$ is a two-sided inverse of $1-\xi\SB$, and any left inverse of
  $1-\xi\SB$ equals it.  This is the workhorse behind properties 1--3.
\end{lemma}

\begin{lemma}[Properties 1--4 of Lemma \texttt{lem\_propTH}]\label{lem:propTH-1234}
  \lean{RBM.Theta_transpose, RBM.Theta_apply_add_right, RBM.Theta_commute_SB,
    RBM.Theta_commute, RBM.Theta_eq_tsum, RBM.sum_Theta_row}
  \uses{def:Theta, lem:Theta-unique, lem:SB-basic}
  $\Th^{(B)}_\xi$ is symmetric and translation invariant, commutes with $\SB$ and
  with $\Th^{(B)}_{\xi'}$, and for $|\xi|<1$ admits the random-walk representation
  $\Th^{(B)}_\xi=\sum_{k\ge0}\xi^k(\SB)^k$.  Moreover
  $\sum_b(\Th^{(B)}_\xi)_{ab}=(1-\xi)^{-1}$.
\end{lemma}

\begin{lemma}[Crude bounds]\label{lem:crude}
  \lean{RBM.norm_Theta_le, RBM.sum_norm_Theta_row_le, RBM.norm_Theta_apply_le}
  \uses{lem:propTH-1234}
  $\max_a\sum_b|(\Th^{(B)}_\xi)_{ab}|\le(1-|\xi|)^{-1}$.  This is \emph{not} the
  sharp bound: the decay length here is governed by $1-|\xi|$, whereas the paper
  needs $\lhat(\xi)=\min(|1-\xi|^{-1/2},L)$.  The two differ enormously when
  $\xi=tm^2$.
\end{lemma}

\chapter{Section 8: proof of Lemma \texorpdfstring{\texttt{lem\_propTH}}{lem-propTH}}

\begin{definition}[Scales, \texttt{(eq\_kappa\_def)}]\label{def:kappa}
  \lean{RBM.kappa, RBM.ellhat}
  $\kappa:=|1-\xi|^{1/2}$ and $\lhat=\lhat(\xi):=\min(\kappa^{-1},L)$.
\end{definition}

\begin{lemma}\label{lem:kappa-basic}
  \lean{RBM.kappa_sq, RBM.kappa_pos, RBM.kappa_mul_ellhat_le_one}\uses{def:kappa}
  $\kappa^2=|1-\xi|$, $\kappa>0$ for $|\xi|<1$, and $\kappa\lhat\le1$.
\end{lemma}

\section{Fourier setup}

\begin{definition}[The symbol, \texttt{(eq\_symbol)}]\label{def:Shat}
  \lean{RBM.Shat}\uses{def:SB}
  For $p\in\TL$,
  $\Shat(p)=\frac{1+2\cos p_1+2\cos p_2}{5}$.
\end{definition}

\begin{lemma}[Plane waves diagonalize $\SB$]\label{lem:eigen}
  \lean{RBM.Shat_eq_cos, RBM.SB_mulVec_char, RBM.norm_Shat_le_one,
    RBM.one_sub_mul_Shat_ne_zero}
  \uses{def:Shat, def:SB}
  $\SB e_p=\Shat(p)e_p$, $|\Shat(p)|\le1$, and $1-\xi\Shat(p)\ne0$ for $|\xi|<1$.
\end{lemma}

\begin{lemma}[Fourier representation, \texttt{(eq\_Fourier\_rep)}]\label{lem:fourier}
  \lean{RBM.Theta_apply_fourier, RBM.inv_mul_sum_chr}
  \uses{lem:eigen, lem:Theta-unique}
  With $x=a-b$,
  \[ (\Th^{(B)}_\xi)_{ab}=K_{\xi,L}(x)=\frac1{L^2}\sum_{p\in\TL}
     \frac{e^{\mathrm{i}p\cdot x}}{1-\xi\Shat(p)} . \]
\end{lemma}

\begin{definition}[\texttt{(eq\_qdef)}]\label{def:q}
  \lean{RBM.qsym}\uses{def:Shat}
  $q(p):=1-\Shat(p)=\frac25\big[(1-\cos p_1)+(1-\cos p_2)\big]$.
\end{definition}

\begin{lemma}[Ellipticity, \texttt{(eq\_elliptic)}]\label{lem:elliptic}
  \lean{RBM.norm_one_sub_mul_Shat_le, RBM.norm_one_sub_mul_Shat_ge,
    RBM.norm_one_sub_mul_Shat_ge_kappa}
  \uses{def:q, def:kappa}
  For $|\xi|<1$,
  \[ \tfrac19\big(\kappa^2+q(p)\big)\;\le\;\big|1-\xi\Shat(p)\big|\;\le\;\kappa^2+q(p). \]
  The paper's proof splits at $\Shat(p)=1/2$; in the low branch its constant
  $1/2$ should read $2/5$, see \texttt{docs/paper-deltas.md}, item 5.
\end{lemma}

\begin{lemma}[\texttt{(eq\_qcomp)}]\label{lem:qcomp}
  \uses{def:q}
  $q(p)\sim|p|_*^2$, where $|p|_*^2=\operatorname{dist}(p_1,2\pi\mathbb Z)^2
  +\operatorname{dist}(p_2,2\pi\mathbb Z)^2$.
  \emph{Work order T2.}
\end{lemma}

\begin{lemma}[Lattice sum]\label{lem:latticesum}
  \uses{lem:qcomp}
  $\sum_{p\in\TL\setminus\{0\}}|p|_*^{-2}\le CL^2\log L$.
  \emph{Work order T3.}
\end{lemma}

\section{Property 5: exponential decay}

\begin{lemma}[$\kappa L<1$]\label{lem:decay-small}
  \uses{lem:fourier, lem:elliptic, lem:latticesum}
  If $\kappa L<1$, so that $\lhat=L$, then separating the zero mode gives
  $|K_{\xi,L}(x)|\le\kappa^{-2}L^{-2}+C\log L\prec\kappa^{-2}L^{-2}$, which is
  \texttt{(prop:ThfadC)} because $|x|_L\le CL=C\lhat$.
  \emph{Work order T4.}
\end{lemma}

\begin{lemma}[Infinite-volume kernel and the contour shift]\label{lem:contour}
  \uses{lem:elliptic}
  $|D_\xi(p+\mathrm{i}\eta)|\ge c(\kappa^2+|p|_*^2)$ for $|\eta|\le c_0\kappa$, and
  hence $|K_{\xi,\infty}(x)|\le C\log(2+\kappa^{-1})e^{-c\kappa|x|}$.
  \emph{Work order T9.}
\end{lemma}

\begin{lemma}[Periodization]\label{lem:periodize}
  \uses{lem:contour, lem:fourier}
  $K_{\xi,L}(x)=\sum_{n\in\mathbb Z^2}K_{\xi,\infty}(x+nL)$, by comparing Fourier
  coefficients rather than by invoking Poisson summation.
  \emph{Work order T10.}
\end{lemma}

\begin{lemma}[Property 5, \texttt{(prop:ThfadC)}]\label{lem:propTH-5}
  \uses{lem:decay-small, lem:periodize}
  $(\Th^{(B)}_\xi)_{ab}\prec e^{-c|a-b|_L/\lhat(\xi)}\big/\big(|1-\xi|\lhat(\xi)^2\big)$.
\end{lemma}

\section{Property 6: derivative bounds}

\begin{lemma}[Case 2: $|s|_L>d/2$]\label{lem:bd-case2}
  \uses{lem:fourier, lem:elliptic, lem:latticesum}
  Both differences are $O(\log L)\prec1$, while $|s|_L/(d+1)\ge c$ and
  $|s|_L^2/(d^2+1)\ge c$.
  \emph{Work order T5.}
\end{lemma}

\begin{lemma}[Dyadic estimate, \texttt{(eq\_dyadic)}]\label{lem:dyadic}
  \uses{lem:elliptic}
  $|\nabla^mK_r(x)|\le C_M\,r^{m+2}(\kappa^2+r^2)^{-1}(1+r|x|_L)^{-M}$, $m=0,1,2$.
  \emph{Work order T11.}
\end{lemma}

\begin{lemma}[Property 6, \texttt{(prop:BD1)}--\texttt{(prop:BD2)}]\label{lem:propTH-6}
  \uses{lem:bd-case2, lem:dyadic, lem:kappa-basic}
  The first- and second-difference bounds of Lemma \texttt{lem\_propTH}.
\end{lemma}

\chapter{Delocalization}

\begin{lemma}[Deterministic core of Theorem \texttt{MR:decol}]\label{lem:deloc}
  \lean{RBM.green, RBM.green_eq_spectral, RBM.im_green_apply_self,
    RBM.sq_norm_eigenvector_le_im_green, RBM.sq_norm_eigenvector_le_of_norm_green_le}
  For Hermitian $H$ with eigenpairs $(\lambda_k,\psi_k)$ and $\eta>0$,
  \[ |\psi_k(x)|^2\;\le\;\sum_l\frac{\eta^2|\psi_l(x)|^2}{(\lambda_k-\lambda_l)^2+\eta^2}
     \;=\;\eta\,\operatorname{Im}G_{xx}(\lambda_k+\mathrm{i}\eta), \]
  so a bound $|G_{xx}|\le C$ gives $|\psi_k(x)|^2\le C\eta$.  The probabilistic
  input (the local law \texttt{MR:locSC}) is a hypothesis here.
\end{lemma}

\chapter{Not formalized}

\begin{remark}\label{rem:random}
  Mathlib has no It\^o calculus for matrix-valued Brownian motion, no matrix SDEs
  and no Dyson Brownian motion.  The loop hierarchy of \S5--\S7 and the inductive
  scheme \texttt{lem:main\_ind} are therefore out of scope; they are to be carried
  as \texttt{axiom} interfaces if and when the logical skeleton needs them.
\end{remark}
